\documentclass[10pt,aspectratio=169]{beamer}
\input{../../_en_preambulo_beamer}% Comandos y macros estadísticas compartidas para las presentaciones Beamer en inglés (Metropolis)

% Conjuntos numéricos básicos
\newcommand{\R}{\mathbb{R}}
\newcommand{\N}{\mathbb{N}}
\newcommand{\Q}{\mathbb{Q}}
\newcommand{\Z}{\mathbb{Z}}
\newcommand{\set}[1]{\left\{ #1 \right\}}
\newcommand{\abs}[1]{\left|#1\right|}
\newcommand{\norm}[1]{\left\| #1 \right\|}

% Comandos estadísticos y probabilísticos del curso
\newcommand{\Var}{\operatorname{Var}}
\newcommand{\cov}{\operatorname{Cov}}
\newcommand{\corr}{\rho}
\newcommand{\comb}[2]{\begin{pmatrix} #1 \\ #2 \end{pmatrix}}
\newcommand{\s}{\sigma}
\newcommand{\particion}{\mathfrak{P}}
\newcommand{\card}[1]{\#\left( #1 \right)}
\newcommand{\seq}[3]{\set{#1}_{#2}^{#3}}
\newcommand{\E}{\mathbb{E}}
\newcommand{\Prob}{\mathbb{P}}

% Entornos pedagógicos y cajas en inglés académico natural (soportando tanto nombres en inglés como en español)
\newenvironment{discussion}{%
  \begin{alertblock}{Discussion Question}%
}{%
  \end{alertblock}%
}
\newenvironment{discusion}{%
  \begin{alertblock}{Discussion Question}%
}{%
  \end{alertblock}%
}

\newenvironment{reflection}{%
  \begin{alertblock}{Classroom Reflection}%
}{%
  \end{alertblock}%
}
\newenvironment{reflexion}{%
  \begin{alertblock}{Classroom Reflection}%
}{%
  \end{alertblock}%
}

\newenvironment{paradox}{%
  \begin{alertblock}{Exploratory Paradox}%
}{%
  \end{alertblock}%
}
\newenvironment{paradoja}{%
  \begin{alertblock}{Exploratory Paradox}%
}{%
  \end{alertblock}%
}

\newenvironment{motivation}{%
  \begin{exampleblock}{Motivation: Data Science in Practice}%
}{%
  \end{exampleblock}%
}
\newenvironment{motivacion}{%
  \begin{exampleblock}{Motivation: Data Science in Practice}%
}{%
  \end{exampleblock}%
}

\newenvironment{quickchallenge}{%
  \begin{block}{Quick Team Challenge}%
}{%
  \end{block}%
}
\newenvironment{retoclase}{%
  \begin{block}{Quick Team Challenge}%
}{%
  \end{block}%
}

\title[$p$-values and decisions]{Section 07.03: $p$-Values and Decisions}\subtitle{Evidence, significance, and interpretation}
\author[J. Castillo Colmenares]{Juliho Castillo Colmenares}\institute{Tecnológico de Monterrey}\date{}
\begin{document}
\begin{frame}[plain]\titlepage\end{frame}
\begin{frame}{Roadmap}\small\begin{enumerate}\item Significance level.\item Definition of the $p$-value.\item Procedure and interpretation.\end{enumerate}\end{frame}
\begin{frame}{Source and scope}\small The notes define $\alpha$, the $p$-value, and a complete decision procedure. This deck follows that path.\begin{block}{Guiding question}How incompatible are the data with $H_0$?\end{block}\end{frame}
\begin{frame}{Significance level}\small $\alpha$ is the maximum tolerated probability of a Type I error. Common values are $0.10$, $0.05$, and $0.01$.\begin{alertblock}{Before observing} Set $\alpha$ before computing the statistic.\end{alertblock}\end{frame}
\begin{frame}{$p$-value}\small\begin{block}{Definition} The $p$-value is the probability of observing a result at least as extreme as the observed result, assuming the null hypothesis is true.\end{block}\end{frame}
\begin{frame}{What it does not mean}\small\begin{itemize}\item It is not the probability that $H_0$ is true.\item It does not measure effect size by itself.\item It depends on the model, statistic, and alternative direction.\end{itemize}\end{frame}
\begin{frame}{Decision rule}\small\[p\leq\alpha\Rightarrow\text{reject }H_0;\qquad p>\alpha\Rightarrow\text{do not reject }H_0.\]\begin{block}{Language} The decision concerns evidence, not logical certainty.\end{block}\end{frame}
\begin{frame}{Alternatives and tails}\small For a one-sided test, compute the $p$-value in one tail; for a two-sided test, use both extremes. The alternative sets the direction.\end{frame}
\begin{frame}{Procedure I}\small\begin{enumerate}\item State $H_0$ and $H_a$.\item Choose $\alpha$.\item Select the statistic and its distribution under $H_0$.\end{enumerate}\end{frame}
\begin{frame}{Procedure II}\small\begin{enumerate}\setcounter{enumi}{3}\item Compute the observed statistic.\item Obtain the $p$-value.\item Compare it with $\alpha$.\end{enumerate}\end{frame}
\begin{frame}{Procedure III}\small\begin{block}{Report} Include the statistic, degrees of freedom, $p$-value, $\alpha$, and a contextual conclusion.\end{block}\end{frame}
\begin{frame}{Application: setup}\small A company claims a mean delivery time of three days. The claim is suspected to be too low; a sample yields a $t$ statistic.\begin{block}{Hypotheses} $H_0:\mu=3$ versus $H_a:\mu>3$.\end{block}\end{frame}
\begin{frame}{Application: statistic}\small With $\bar x$, $s$, and $n$,\[t=\frac{\bar x-3}{s/\sqrt n}.\]The degrees of freedom are $n-1$.\end{frame}
\begin{frame}{Application: $p$-value}\small The $p$-value is the right-tail area beyond the observed statistic under $t_{n-1}$.\end{frame}
\begin{frame}{Application: decision}\small If $p\leq0.05$, reject $H_0$ and report evidence of a larger mean. If $p>0.05$, evidence is insufficient.\end{frame}
\begin{frame}{Interpretation I}\small Statistical significance does not imply practical relevance. Pair the $p$-value with an interval and an effect measure.\end{frame}
\begin{frame}{Interpretation II}\small A small $p$-value indicates incompatibility with $H_0$ under the model; it does not prove the alternative or practical importance.\end{frame}
\begin{frame}{Important considerations}\small\begin{itemize}\item Choose $\alpha$ before analysis.\item Do not change the alternative after seeing data.\item Avoid unadjusted repeated queries.\end{itemize}\end{frame}
\begin{frame}{Reporting application}\small\begin{block}{Template} ``The statistic was $T=\cdots$, with $p=\cdots$. At $\alpha=\cdots$, we [reject/do not reject] $H_0$.''\end{block}\end{frame}
\begin{frame}{Procedure summary}\small\begin{enumerate}\item Hypotheses.\item Level.\item Statistic.\item $p$-value.\item Decision.\item Interpretation.\end{enumerate}\end{frame}
\begin{frame}{Practical conclusion}\small A reproducible decision preserves the alternative, selected level, and $p$-value so another reader can review both calculation and interpretation.\end{frame}
\begin{frame}{Chapter connection}\small $p$-values complement confidence intervals and prepare tests for means, proportions, and variances.\end{frame}
\end{document}
